
\section{Neutrino oscillation fitting method}
\label{sec:neutrino_oscillation_fitting}



%--------------------------------------------------------------------------------------------------------------------
\subsection {Fit method}
\label{sec:oscfit_method}

Several choices exist for the oscillation parameters that could be fitted. 
As it has been already introduced in Sec. \ref{sec:osc_params_choice}, this analysis is 
fitting for $\sin^2(2\theta_{23})$ in order to take into account the effect of the octant 
of this parameter (as it was already discussed, during the review process of the T2K Run 1+2+3 $\nu_{\mu}$-disappearance PRL \cite{T2KRun123Disapp} it was found that the choice of the octant had a large effect on the confidence contour in $\sin^2(2\theta_{23})$).

There are also different choices for the mass-squared splitting parameter: typically $|\Delta m^{2}_{32}|$ is used in $\nu_{\mu}$-disappearance 
analyses, but, as explained in Sec.\ref{sec:osc_params_choice}, the physical meaning of this parameter is different depending on the mass hierarchy, since it is the largest mass-squared splitting for the inverted mass
hierarchy and the second largest for the normal mass hierarchy. 
It is therefore not correct to compare results obtained under the different mass hierarchy hypotheses.
For this reason, this analysis is fitting for the mass-squared difference defined as
$\Delta m^{2} = \left(m_{3}^{2} - \frac {m_{1}^{2}+m_{2}^{2}} {2} \right)$,
whose absolute value is the same for both mass hierarchies  \cite{FogliLisi}.
Another option, used in the {\em VALOR} Run 1+2+3+4 $\nu_{\mu}$-disappearance analysis \cite{OurRun1234DisappAnalysis3Flavour}, is to fit $|\Delta m^{2}_{32}|$ for the normal hierarchy and $|\Delta m^{2}_{31}|$ for the inverted hierarchy. 

The {\em VALOR} joint oscillation analysis, like the {\em VALOR} disappearance analysis, uses a binned likelihood-ratio method.
Measurements of the oscillation parameters $\sin^{2}\theta_{23}$, $|\Delta m^{2}|$,  $\sin^{2}\theta_{13}$ and $\delta_{CP}$\footnote{
   As already mentioned, a study has been made in which $\sin^{2}\theta_{12}$ and $\delta m^{2}$ were allowed to float in the fit and the results are shown in Appendix \ref{app:solar}. 
   The effect was negligible and it was decided to keep them fixed to the values in Tab.\ref{tab:solarparameters}.
}
are obtained by comparing the observed and predicted SuperK reconstructed energy spectra
for single $\mu$-like ring events and single e-like ring events. 

Let $N_{\mu}$ and $N_{e}$ be the number of reconstructed energy bins
considered for the single $\mu$-like ring and single e-like ring samples respectively, 
and ${\bf a}$ be a $1 \times N_{s}$ - dimensional array of systematic (nuisance) 
parameters that affect the SuperK reconstructed energy spectrum prediction of either sample 
(e.g. parameters that control the beam, neutrino interaction simulations and 
the detector response).
Best-fit values are obtained by minimizing: 

\begin{equation}
\begin {split}
\displaystyle
%\chi^{2}   &=                                                                                                                                                                                       
 -2 \, \textrm{ln} \lambda (|\Delta m^{2}|, sin^{2}\theta_{23}, \sin^{2}\theta_{13}, \delta_{CP}; {\bf a}) &=
  2 \cdot \sum_{r=0}^{N_{\mu}-1}
   \left[
    (N^{exp}_{SK;\mu,r}-N^{obs}_{SK;\mu,r}) + 
     N^{obs}_{SK;\mu,r} \cdot ln\left(\frac{N^{obs}_{SK;\mu,r}}{N^{exp}_{SK;\mu,r}}\right)  
   \right]  \\
    \displaystyle
    &+ 2 \cdot \sum_{r=0}^{N_{e}-1}
    \left[
    (N^{exp}_{SK;e,r}-N^{obs}_{SK;e,r}) + 
     N^{obs}_{SK;e,r} \cdot ln\left(\frac{N^{obs}_{SK;e,r}}{N^{exp}_{SK;e,r}}\right)  
   \right]  \\
  &+  ({\bf a}-{\bf a_{0}})^{T} \cdot {\bf C}^{-1} \cdot ({\bf a}-{\bf a_{0}})
\label{eq:chisquare}
\end{split}
\end{equation}
where
$N^{obs}_{SK;\mu,r}$ ($N^{obs}_{SK;e,r}$) is the observed number of single $\mu$-like ring (single e-like ring) events in the $r^{th}$ bin, 
and $N^{exp}_{SK;\mu,r}$ ($N^{exp}_{SK;e,r}$)
is the corresponding expected number of events; this expected number is a function of 
$|\Delta m^{2}|$, $\sin^{2}\theta_{23}$, $\sin^{2}\theta_{13}$ and $\delta_{CP}$
and the nuisance parameter vector ${\bf a}$.
Also ${\bf a_{0}}$ is a $1 \times N_{s}$ - dimensional array containing the nominal values of
the systematic parameters, ${\bf a}^{T}$ is the transpose of ${\bf a}$,
and ${\bf C}$ is the systematic parameter covariance matrix of dimension $N_{s} \times N_{s}$.
For the stand-alone $\nu_{\mu}$-disappearance ($\nu_{e}$-appearance) fit, 
only the single $\mu$-like ring (single e-like ring) sample is used and the 
corresponding term in Eq.\ \ref{eq:chisquare} is dropped.

The likelihood-ratio method used in this analysis is equivalent to the extended maximum likelihood
method used in the alternative T2K oscillation analysis 
(see \cite{OurRun1DisappAnalysis,StatAnalysisCowan} for more details).
The advantage of the likelihood ratio method, compared with the extended maximum-likelihood method, is that
in the large-sample limit, the quantity $-2 \: \textrm{ln} \: \lambda (\sin^{2}\theta_{23},|\Delta m^{2}_{32}| \; \textrm{or} \; |\Delta m^{2}_{31}|; {\bf a})$
in Eq.\ref{eq:chisquare} has a $\chi^{2}$ distribution and it can therefore be used as a goodness-of-fit test.
In this note,
$-2 \: \textrm{ln} \: \lambda (\sin^{2}\theta_{23},|\Delta m^{2}_{32}| \; \textrm{or} \; |\Delta m^{2}_{31}|; {\bf a})$ and
$\chi^{2}(\sin^{2}\theta_{23},|\Delta m^{2}_{32}| \; \textrm{or} \; |\Delta m^{2}_{31}|; {\bf a})$ are used interchangeably.

The minimization is performed with Minuit \cite{Minuit}, using the Migrad method.
If Migrad does not converge, the starting values of the oscillation parameters are re-initialised:
the new initial values are randomly drawn from uniform distributions in the ranges 
$2 \times 10^{-3}$ eV$^{2}$/c$^{4}$ $\leq$ $|\Delta m^{2}|$ $\leq$ $3 \times 10^{-3}$  eV$^{2}$/c$^{4}$,
$0.015 \leq \sin^{2}\theta_{13} \leq 0.035 $ (or $0.05 \leq \sin^{2}2\theta_{13} \leq 0.15 $), 
$-\pi \leq \delta_{CP} \leq \pi$ and
either 0.35 $\leq$ $\sin^{2}\theta_{23}$ $\leq$ 0.50 or 0.53 $\leq$ $\sin^{2}\theta_{23}$ $\leq$ 0.68 
(the new initial value of $\sin^{2}\theta_{23}$ is in the same octant as the previous initial value).  
The fit is retried with the new initial values, and this is repeated up to 25 times.
If the fit still fails to converge after 25 attempts, the Simplex algorithm is used.

The search for the minimum value of 
$-2 \, \textrm{ln} \lambda(|\Delta m^{2}|, \sin^{2}\theta_{23}, \sin^{2}\theta_{13}, \delta_{CP}; {\bf a})$
is performed in the range 
1.0 $\times$10$^{-3}$ eV$^{2}$/c$^{4}$ $\leq$ $|\Delta m^{2}|$ $\leq$ 6.0 $\times$10$^{-3}$ eV$^{2}$/c$^{4}$,
0.0 $\leq$ $\sin^{2}\theta_{23}$ $\leq$ 1.0,
0.0 $\leq$ $\sin^{2}\theta_{13}$ $\leq$ 1.0 and
$-\pi$ $\leq$ $\delta_{CP}$ $\leq$ $\pi$.
After a successful fit, The Hesse algorithm is called to improve Minuit's estimates of the
errors on the fitted parameters.

In order to find a solution in each octant for $\sin^{2}\theta_{23}$, and since it is unlikely that Minuit jumps to the minimum
in the octant mirror to the one it was initilized in, each fit is repeated twice:
in the first fit, the initial value is $\sin^{2}\theta_{23} = 0.6$, while the second fit uses as initial value the mirror point in the other octant of the best fit of $\sin^{2}\theta_{23}$ from the first fit (taking into account the best-fit value of  $\sin^{2}\theta_{13}$ from the first fit to calculate the correct mirror point: $(\sin^{2}\theta_{23})_{mirror}$ = center $\pm |$ center - $(\sin^{2}\theta_{23})_{bf}|$ where the center of symmetry (value of maximal disappearance) is calulated as
center = $(2(1-(\sin^{2}\theta_{13})_{bf}))^{-1}$ as explained in Appendix E of \cite{OurJointAnalysis}). 
An simpler method, setting initial values $\sin^{2}\theta_{23} = 0.6$ and $\sin^{2}\theta_{23} = 0.4$ for the two fits,
 was also tested, resulting as well in a solution in each octant as expected.

All the 64 systematic parameters are allowed to float in the fit, and are restricted to the range [-5$\sigma_{s}$, +5$\sigma_{s}$]
where $\sigma_{s}$ is the one standard deviation error assigned to each systematic parameter\footnote{
 Such an error is defined for all but one of the systematic parameters considered in this analysis. 
 The exception is the systematic parameters $f_{SF}$ which parameterizes the uncertainty on nuclear modelling,
 switching between the Relativistic Fermi Gas (RFG) and Spectral Function (SF) models. The allowed range of the parameter $f_{SF}$ 
 in the oscillation fit is [0,1].}.
Values of systematic parameters that give a negative predicted number of events in any reconstructed energy bin are not allowed. 
If this scenario arises, the number of events is changed to +1E-8 in that bin. This is done for the sum of all the interaction 
modes rather than for each mode separately.

%..................................
%

\subsection{Adding reactor constraint}
\label{sec:reactor_constraint}

In this technical note, the results of the Run 1+2+3+4 dataset  will be presented in two ways: one was done fitting 
only T2K dataset and one will be including the reactor constraint in the fit.
This constraint works adding a penalty term to the $\chi^2$ function, proportional to the difference
between the best-fit value of $\sin^{2}\theta_{13}$ and the one used as constraint. 
The $\chi^2_{reactor}$ is calculated following Eq. \ref{eq:reactor}, where the values of 
$(\sin^{2}\theta_{13})_{reactor}$ and $\sigma_{reactor}$ used were the latest 2013 PDG values, 
coming from a weighted average of the results from the three reactor experiments Daya Bay, RENO and Double Chooz \cite{PDG2013}:
$(\sin^{2}\theta_{13})_{reactor}$ = 0.095 and $\sigma_{reactor}$ = 0.01.

\begin{equation}
\chi^2_{reactor} = (\frac{\sin^{2}\theta_{13} - (\sin^{2}\theta_{13})_{reactor} }{\sigma_{reactor}})^2
\label{eq:reactor}
\end{equation}

%..................................
%

\subsection{Goodness-of-fit tests}
\label{sec:goodness_of_fit}

In this analysis, the procedure to be used for computing the goodness-of-fit will be the same one used in the {\em VALOR} Run 1+2+3 joint oscillation analysis \cite{OurJointAnalysis}. {\color{red}{Note: goodness-of-fit tests have not yet been calculated, they will be included in a later version of this technical note.}}

The problem of defining a general goodness-of-fit for sparsely populated distributions
was discussed in detail during the development of the Run 1+2+3 stand-alone $\nu_\mu$-disappearance 
and $\nu_{e}$-appearance analyses. See for example \cite{KarlenPvalue}.
It was pointed out that, for the Run 1+2+3 statistics, the shape information is not
as helpful in forming p-values as the rate information.
It was decided that  p-values will be calculated using the rate and the reconstructed energy distribution,
and that the distribution will use a very coarse binning
scheme to ensure that there are sufficient numbers of events in each bin.

A similar approach is the one used for the joint 3-flavour analysis presented in \cite{OurJointAnalysis}.
For the single $\mu$-like ring sample, 5 reconstructed energy bins were used:
0-0.4, 0.4-0.7, 0.7-1.0, 1.0-2.0 and 2.0-30.0 GeV.
For the single e-like ring sample, 3 reconstructed energy bins were used:
0-0.35, 0.35-0.80, 0.80-1.25 GeV.
In this note, p-values for the joint 3-flavour fit are computed following 
the 5+3 binning scheme described above. When p-values were quoted for the stand-alone
$\nu_{\mu}$-disappearance ($\nu_{e}$-appearance) analysis, the corresponding 
5 (3) bin scheme of the single $\mu$-like ring (single e-like ring) event sample was used.

The goodness-of-fit tests that will be shown in this technical note are built using 1000 toy MC experiments
created at the best-fit point to be tested, with statistical and systematic fluctuations.
Then, each toy MC experiment is fitted in the coarse binning described above, 
including the four oscillation parameters and all systematics in the joint 3-flavour fit.


%..................................
%

\subsection{Construction of confidence regions}

In this technical note, confidence regions are constructed using the constant $\Delta\chi^2$ method \cite{PDG2012} (pp. 390-401).
68\% and 90\% CL limits are reported for each of $|\Delta m^{2}|$, $sin^{2}\theta_{23}$, $sin^{2}\theta_{13}$ and $\delta_{CP}$
oscillation parameters. In addition, 2-dimensional confidence regions for 68\% and 90\% CL will be shown for the following combinations of parameters: 
($sin^{2}\theta_{23}$, $|\Delta m^{2}|$),
($sin^{2}\theta_{13}$, $|\Delta m^{2}|$),
($sin^{2}\theta_{13}$, $\delta_{CP}$) and
($sin^{2}\theta_{23}$, $sin^{2}\theta_{13}$).
When introducing the reactor constraint, the following additional 2-dimensional confidence regions are also presented:
 ($|\Delta m^{2}|$,$\delta_{CP}$) and
($sin^{2}\theta_{23}$,$\delta_{CP}$).


The first step to construct confidence contours in an oscillation parameter space 
is to build the $\Delta\chi^2$ surface in that space.
The oscillation parameter space is divided into a regular grid.
The grid limits and spacing was used in all confidence regions is presented in Tab. \ref{tab:grid_binning}.

\begin{table}[htbp]
\center
\begin{tabular}{|p{3.5cm}|p{1.5cm}|p{1.5cm}|p{1.5cm}|p{1.cm}|}
\hline 
Parameter & Min & Max & Step & Bins \\
\hline \hline
$sin^{2}\theta_{23}$ & 0.3 & 0.7 & 0.01 & 41 \\
\hline
$|\Delta m^{2}|$ ($10^{-3} eV^2/c^4$) & 2.0 & 3.0 & 0.05 & 21 \\
\hline
$sin^{2}\theta_{13}$ & 0.0 & 0.1 & 0.005 & 21 \\
\hline
$sin^{2}2\theta_{13}$ & 0.0 & 0.4 & 0.02 & 21 \\
\hline
$sin^{2}\theta_{13}$ (reactor) & 0.015 & 0.035 & 0.001 & 21 \\
\hline
$sin^{2}2\theta_{13}$ (reactor) & 0.06 & 0.14 & 0.004 & 21 \\
\hline
$\delta_{CP}$ & $-\pi$ & $\pi$ & $2\pi/50$ & 51 \\
\hline
\end{tabular}
\caption{
  Grid limits and spacing used in all confidence regions. 
}
\label{tab:grid_binning}
\end{table}

For the joint 3-flavour analysis confidence regions,
the $\chi^2$ is minimized at each grid point with respect to all 64 nuisance parameters considered in this analysis
and all oscillation parameters not shown on the plot (and not considered to be fixed).

The confidence regions for the joint analysis are constructed by finding the global best-fit point
(fitting all oscillation parameters and all systematic parameters at the same time)
and calculating the lines of {\em constant} $\Delta{\chi^{2}}$ from this point\footnote{
For 1 fit parameter, the critical values of $\Delta{\chi^{2}}$ used are 
1.00 (68\% CL), 2.71 (90\% CL), 3.84 (95\% CL), 6.63 (99\% CL).
For 2 parameters fitted simultaneously, the critical values of $\Delta{\chi^{2}}$ used are 
2.30 (68\% CL), 4.61 (90\% CL), 5.99 (95\% CL), 9.21 (99\% CL).
}. A special procedure is followed 
to construct the confidence regions in ($sin^{2}\theta_{13}$, $\delta_{CP}$) space, called \textit{raster scan}, 
where a best-fit value of  $sin^{2}\theta_{13}$ is found for each fixed value of $\delta_{CP}$, resulting in a line
 of best-fit points instead of a single best-fit point. Then, confidence intervals in $sin^{2}\theta_{13}$
 at a fixed $\delta_{CP}$ are found with respect to each best-fit $sin^{2}\theta_{13}$ using the 1D
 $\Delta{\chi^{2}}$ critical values \cite{FeldmanCousins}.

The method of Feldman and Cousins could be used to find the critical values of $\Delta{\chi^{2}}$
for the confidence regions of the joint oscillation analysis.
In that case, the prescription to find the Feldman-Cousins confidence regions would be imitating the
one used for constructing 1D confidence regions in the {\em VALOR} $\nu_{\mu}$ disappearance analysis
\cite{OurRun1234DisappAnalysis3Flavour}. Independent Feldman-Cousins confidence regions would be
computed for each one of the 2-dimensional parameter spaces separately, using the results in this 
analysis to select the true input values of the other two oscillation parameters which are fitted for each
2-dimensional parameter space, following for them the $\Delta\chi^2$ surface obtained from the
Run 1+2+3+4 dataset fit obtained in the parameter space defined by those other two oscillation parameters.
Using these $\Delta\chi^2$ surfaces, the other two oscillation parameters would be varied following the
preference of T2K Run 1+2+3+4 data. 

For the inclusion of systematics in the Feldman and Cousins method, the same procedure as used for the 
{\em VALOR} $\nu_{\mu}$ disappearance analysis would be followed. 
To construct a 2-dimensional confidence region, for instance (sin$^{2}\theta_{23}$, $|\Delta m^{2}$),
where N systematic parameters are included in the p.d.f. and all of these N systematics are allowed to float in the fit,
the value of $\Delta\chi^{2}_{critical}$ at each true (sin$^{2}\theta_{23}$, $|\Delta m^{2}|$) grid point 
would be calculated from an ensemble of toy MC experiments,
each one of which is generated with the true oscillation parameters of that point and with all N systematics randomised.
Then, for each toy MC experiment, both sin$^{2}\theta_{23}$, $|\Delta m^{2}|$ and all N nuisance parameters along with 
sin$^{2}\theta_{13}$ and $\delta_{CP}$ are fitted
simultaneously to obtain the value of $\chi^{2} \textrm{(MC, best fit)}$ used in Eq. \ref{eq:deltachi2mc}.
A separate fit, with sin$^{2}\theta_{23}$ and $|\Delta m^{2}|$ fixed to their true values
and with all N nuisance parameters allowed to float along with sin$^{2}\theta_{13}$ and $\delta_{CP}$, 
is performed to obtain the value of  $\chi^{2} \textrm{(MC, true)}$ (i.e.\ $\chi^{2}$ is \emph{independently} minimised, for \emph{each} toy MC experiment, with respect to the nuisance parameters).

Feldman-Cousins values of critical $\Delta\chi^2$ vs $\delta_{CP}$ have already been calculated and presented in Sec. \ref{sec:results:fc_deltacp}. The procedure used to calculate them is explained in that section. 
